\begin{teilaufgaben}
\item
The function $u(x)$ is defined on the interval $\Omega = [-1,1]$ and
satisfies in $\Omega$ 
\[
u''(x) + 2 \cdot u'(x) = 2
\]
and $u(-1) = 0, u(1) = 0$ on the boundary of $\Omega$. 

Determine approximate values for $u(-1/2)$, $u(0)$ and $u(1/2)$. 
Use centered finite differences with $h = 1/2$. 
\item
The function $u(x,y)$ is defined on the rectangle
$\Omega = [-1, 1] \times [0,3]$ and satisfies in $\Omega$
\[
\Delta u(x,y) = e^x \cdot y 
\]
and $u(x,y) = 0$ on the boundary of $\Omega$.

Determine approximate values for $u(0,1)$ and $u(0,2)$.
Use finite differences with $h = 1.$
\end{teilaufgaben}

\begin{loesung}
\begin{teilaufgaben}
\item
By the method of finite differences we have
\begin{equation*}
\renewcommand{\arraycolsep}{2pt}
\renewcommand{\arraystretch}{1.3}
\begin{array}{rclclclcrclclcl}
4&\cdot&\big(\tilde u(-1) &-& 2 \cdot \tilde u(-1/2) &+& \tilde u(0)
&\llap{$\big)$}\hspace{2pt}+&
2&\cdot&\big(\tilde u(0) &-& \tilde u(-1)
&\llap{$\big)$}\hspace*{2pt}=&
2,
\\
4&\cdot&\big(\tilde u(-1/2) &-& 2 \cdot \tilde u(0) &+& \tilde u(1/2)
&\llap{$\big)$}\hspace{2pt}+&
2&\cdot&\big(\tilde u(1/2) &-& \tilde u(-1/2)
&\llap{$\big)$}\hspace*{2pt}=&
2,
\\
4&\cdot&\big(\tilde u(0) &-& 2 \cdot \tilde u(1/2) &+& \tilde u(1)
&\llap{$\big)$}\hspace{2pt}+&
2&\cdot&\big(\tilde u(1) &-& \tilde u(0)
&\llap{$\big)$}\hspace*{2pt}=&
2.
\end{array}
\end{equation*}
This leads to a system of three linear equations:
\begin{equation}
\begin{pmatrix*}[r] -8 & 6 & 0 \\  2 & -8 & 6 \\  0 & 2 & -8 \end{pmatrix*}
\cdot
\begin{pmatrix} \tilde u(-1/2)  \\   \tilde u(0) \\ \tilde u(1/2) \end{pmatrix}
=
\begin{pmatrix*}[r] 2  \\  2 \\  2 \end{pmatrix*}.
\tag{\bf{2P}}
\end{equation}
Its solution is
\begin{equation}
\big(\tilde u(-1/2),  \tilde u(0), \tilde u(1/2)\big) = (-0.85 , -0.8, -0.45).
\tag{\bf{1P}}
\end{equation}

\item
By the five-point-star operator we have
\[
\tilde u(0,2) - 4 \cdot \tilde u(0,1) = 1,
\]
\[
\tilde u(0,1) - 4 \cdot \tilde u(0,2) = 2.
\]
This leads to a system of two linear equations:
\begin{equation}
\begin{pmatrix*}[r] -4 & 1 \\  1 & -4 \end{pmatrix*}
\cdot
\begin{pmatrix} \tilde u(0,1)  \\   \tilde u(0,2)  \end{pmatrix}
=
\begin{pmatrix*}[r] 1  \\  2 \end{pmatrix*}.
\tag{\bf{2P}}
\end{equation}
Hence
\begin{equation}
\tilde u(0,1) = -0.4 \qquad\text{and}\qquad \tilde u(0,2) = -0.6.
\tag{\bf{1P}}
\end{equation}
\end{teilaufgaben}
\end{loesung}
